\documentclass[11pt, a4paper]{article}

% --- Packages ---
\usepackage{fontspec}
\setmainfont{Calibri}[
    Path = /Applications/Microsoft Word.app/Contents/Resources/DFonts/,
    Extension = .ttf,
    UprightFont = Calibri,
    BoldFont = Calibrib,
    ItalicFont = Calibrii,
    BoldItalicFont = Calibriz
]
\usepackage[margin=2.5cm]{geometry}
\usepackage{setspace}
\setstretch{1.5}
\usepackage{amsmath}
\usepackage{amsfonts}
\usepackage{amssymb}
\usepackage{titlesec}
\usepackage{graphicx}
\usepackage{booktabs}
\usepackage{tabularx}
\usepackage{longtable}
\usepackage{float}
\usepackage{caption}
\usepackage{subcaption}
\usepackage[style=apa, backend=biber, sorting=nyt]{biblatex}
\addbibresource{references.bib}
\usepackage{csquotes}
\usepackage{enumitem}
\usepackage{multirow}
\usepackage{adjustbox}
\usepackage{pdflscape}
\usepackage{siunitx}

% --- Hyperlinks ---
\usepackage{hyperref}
\hypersetup{
    colorlinks=true,
    linkcolor=black,
    filecolor=magenta,
    urlcolor=blue,
    citecolor=black,
    pdftitle={Master Thesis -- Passive Investing and Active Alpha},
    pdfpagemode=UseOutlines,
}

% --- Section Formatting ---
\titleformat{\section}{\large\bfseries}{\thesection}{1em}{}
\titleformat{\subsection}{\normalsize\bfseries}{\thesubsection}{1em}{}
\titleformat{\subsubsection}{\normalsize\itshape}{\thesubsubsection}{1em}{}
\setlength{\parindent}{0pt}
\setlength{\parskip}{0.8em}

% --- Custom commands ---
\newcommand{\passiveshare}{\textit{PassiveShare}}

\begin{document}

% ============================================================================
% TITLE PAGE
% ============================================================================
\begin{titlepage}
    \centering
    \vspace*{2cm}
    \textsc{\Large BI Norwegian Business School}\\[1.5cm]
    \textsc{\large Master Thesis}\\[0.5cm]

    \rule{\linewidth}{0.5mm} \\[0.4cm]
    {\huge \bfseries The Association Between Passive Investing and Active Management Alpha: Evidence from U.S.\ Mutual Funds} \\[0.4cm]
    \rule{\linewidth}{0.5mm} \\[1.5cm]

    \begin{minipage}{0.8\textwidth}
        \begin{flushleft} \large
            \textbf{Study Programme:} MSc in Finance \\
            \textbf{Supervisor:} Giovanni Pagliardi \\[0.5cm]
            \textbf{Authors:} \\
            Bendic Alexander Utvaag \\
            H\r{a}kon Julius St{\o}rholt \\
        \end{flushleft}
    \end{minipage}
    \vfill
    \begin{center}
        BI Norwegian Business School \\
        Spring 2026
    \end{center}
\end{titlepage}

\newpage

% ============================================================================
% ABSTRACT
% ============================================================================
\begin{abstract}
\noindent
This thesis examines the association between the growth of passive investing and the aggregate, asset-weighted alpha of active mutual fund managers in the United States. Using a survivorship-bias-free sample of approximately 29,800 active and 3,700 passive U.S.\ mutual funds from Morningstar Direct (December 1996 to January 2025), we estimate fund-level risk-adjusted alpha under the CAPM, Fama--French three-factor, and Fama--French five-factor models. Due to a 36-month rolling estimation window, the regression sample begins in December 1999, yielding 302 monthly observations.

We document three findings. First, we demonstrate that a specification we initially employed, regressing net alpha on passive share with fee-spread controls, produces a spurious positive relationship due to the non-stationarity and mechanical collinearity of the fee-spread variable. This methodological pitfall may also tempt future researchers working with similar data. Second, when the specification is corrected using gross alpha as the dependent variable, Newey--West HAC standard errors, and the exclusion of fee spread, the association between passive share and active alpha is weakly negative under the CAPM ($\beta = -0.620$, $p = 0.035$) and statistically indistinguishable from zero under the Fama--French three-factor and five-factor models. We designate the five-factor model as our primary specification and consider a finding robust only if it achieves significance in at least two of three factor models; by this criterion, the main effect of passive share is not robust. Third, an exploratory analysis of stress-conditional interactions reveals that during periods of elevated credit spreads and high passive concentration, multifactor alpha declines significantly. This is the opposite of the Grossman--Stiglitz prediction and is consistent with crowding or liquidity-withdrawal channels, though we emphasize that these exploratory results do not survive multiple-testing correction and should be treated as hypothesis-generating.

\end{abstract}

\newpage
\tableofcontents
\newpage
\listoftables
\listoffigures
\newpage

% ============================================================================
% 1. INTRODUCTION
% ============================================================================
\section{Introduction}
\label{sec:introduction}

\subsection{Background and Motivation}

Over the past three decades, the United States asset management industry has undergone a structural transformation. Capital has migrated from actively managed mutual funds, where portfolio managers select securities on the basis of fundamental analysis, toward passive index funds and exchange-traded funds (ETFs) that mechanically replicate a benchmark. In our sample, the share of total mutual fund assets under management (AUM) held in passive vehicles grew from approximately 5.8\% in December 1996 to over 51\% by January 2025, with passive funds surpassing active funds in total AUM for the first time in August 2024.

This shift has been driven by two reinforcing forces. First, a large body of academic research has shown that the average actively managed fund fails to outperform its benchmark after fees \parencite{fama2010, carhart1997, barras2010}. Second, the emergence of low-cost index products from providers such as Vanguard, BlackRock, and State Street has created a compelling cost alternative, with asset-weighted passive fund fees declining from approximately 0.25\% per year in the late 1990s to below 0.10\% by 2024 in our data.

This migration raises a fundamental question for financial economics: what happens to market efficiency when an increasing share of capital is managed by investors who do not engage in price discovery? The conventional view, rooted in the Efficient Market Hypothesis \parencite{fama1970}, holds that prices reflect available information because active traders compete to exploit mispricings. The Grossman--Stiglitz paradox \parencite{grossman1980} formalizes this intuition: if prices fully reflected all information, no investor would have an incentive to bear the cost of information gathering, but if no one gathers information, prices cannot be informative. By extension, if the fraction of informed traders declines as capital moves to passive strategies, the remaining active managers should face a less competitive environment, theoretically increasing their ability to generate risk-adjusted excess returns (alpha).

Yet the empirical evidence presents a more conflicting picture. Despite the massive growth of passive investing, aggregate measures of active fund performance have not improved. If anything, active managers as a group have continued to underperform their benchmarks after fees \parencite{fama2010}. This disconnect between the theoretical prediction and the observed evidence motivates our investigation.

\subsection{Research Question}

This thesis seeks to answer the following research question:

\begin{quote}
    \textit{What is the empirical association between the rising market share of passive investment and the aggregate, asset-weighted alpha of active fund managers in the U.S.\ mutual fund market over the period 1999 to 2025?}
\end{quote}

We emphasize that our research design identifies an association, not a causal effect. We lack a credible instrument for passive growth and cannot exploit a natural experiment; our findings therefore describe the observed co-movement of passive share and alpha, not the impact of one on the other.

\subsection{Contributions}

Our study makes three contributions. First, we provide a methodological contribution by documenting a specification pitfall: regressing net alpha on passive share with fee-spread controls produces a spurious positive coefficient. We discovered this artifact in our own initial analysis and subsequently identified the mechanism. Fee spread is both non-stationary (Augmented Dickey--Fuller $p = 0.97$; KPSS $p < 0.01$) and mechanically collinear with passive share (both trend monotonically over time), so conditioning on fee spread reverses the sign of the passive share coefficient. The near-identical magnitude of the sign flip in both gross and net alpha specifications (Panel C of our results) confirms that the distortion arises from collinearity rather than the mechanical fee link in the dependent variable. We frame this as a pitfall that we initially committed and that may tempt future researchers working with similar data.

Second, we contribute to the empirical literature on the Grossman--Stiglitz hypothesis by documenting that, once this methodological artifact is removed, the association between passive share and active gross alpha is weakly negative under the CAPM and statistically indistinguishable from zero under the Fama--French three-factor and five-factor models. We designate FF5 as our primary specification and require significance in at least two of three factor models for a finding to be considered robust. By this criterion, we do not find robust evidence of a linear association between passive share and aggregate alpha. This null result is, however, subject to important power limitations: with only 302 monthly observations, our test has limited ability to detect economically small effects, particularly in the multifactor models.

Third, we provide exploratory evidence that the association between passive share and alpha may be conditional on market regime. During periods of elevated credit spreads and high passive concentration, multifactor alpha declines significantly. This is the opposite of the Grossman--Stiglitz prediction and is directionally consistent with crowding or liquidity-withdrawal channels. We emphasize that this finding emerges from a single factor model (FF3), does not survive formal multiple-testing correction \parencite{benjamini1995}, and should be treated as hypothesis-generating rather than confirmatory.

\subsection{Structure}

The remainder of this thesis is organized as follows. Section~\ref{sec:literature} reviews the theoretical foundations and empirical literature. Section~\ref{sec:methodology} describes the data, variable construction, and econometric approach. Section~\ref{sec:results} presents the empirical results across four panels of specifications. Section~\ref{sec:discussion} discusses the findings and their implications. Section~\ref{sec:conclusion} concludes.


% ============================================================================
% 2. LITERATURE REVIEW
% ============================================================================
\section{Literature Review}
\label{sec:literature}

\subsection{Theoretical Foundations}

\subsubsection{The Efficient Market Hypothesis}

The theoretical starting point for our analysis is the Efficient Market Hypothesis (EMH), as formalized by \textcite{fama1970}. In its semi-strong form, the EMH states that asset prices fully reflect all publicly available information, making it impossible for investors to consistently earn abnormal returns through fundamental analysis. The strong form extends this to include private information. Under the EMH, the returns to active management should, on average, equal zero before costs and be negative after costs.

\subsubsection{The Grossman--Stiglitz Paradox}

\textcite{grossman1980} challenged the logical foundations of the EMH by demonstrating that perfectly efficient markets are internally inconsistent. Their model shows that if prices fully reflected all information, no investor would have an incentive to bear the cost of acquiring that information. But if no one acquires information, prices cannot reflect it. This paradox implies that equilibrium must feature some degree of informational inefficiency, sufficient to compensate informed traders for their costs.

The Grossman--Stiglitz framework yields a prediction for our setting: as capital moves to passive funds (which do not acquire or trade on information), the returns to informed trading should rise, potentially drawing capital back toward active management until a new equilibrium is reached. In formal terms, the equilibrium level of informational efficiency is determined by the ratio of informed to uninformed traders, and a decline in this ratio should produce a measurable increase in the alpha available to the remaining informed participants.

\textcite{pedersen2015} extends this logic in his ``efficiently inefficient'' framework, arguing that markets are efficient enough to discourage most investors from paying for active management, yet inefficient enough to compensate the marginal informed trader. Under this view, passive investing shifts the equilibrium toward greater inefficiency, but only gradually and within bounds.

\subsubsection{The Arithmetic of Active Management}
\label{subsub:sharpe_arithmetic}

\textcite{sharpe1991} provides an accounting identity that constrains the aggregate performance of active managers. Since active and passive investors together constitute the entire market, and passive investors by definition earn the market return, the dollar-weighted gross return of all active investors must also equal the market return. After costs, the average active investor must therefore underperform. This identity does not preclude a subset of skilled managers from generating positive alpha, but it establishes that active management is a zero-sum game in aggregate before costs, and a negative-sum game after costs.

This constraint has important implications for our empirical strategy. Because the Sharpe arithmetic pins the AUM-weighted average gross return of active managers to approximately the market return, the aggregate gross alpha that we measure should fluctuate around zero regardless of the level of passive share. Our measured alpha can deviate from zero for three reasons. First, we estimate benchmark-relative CAPM alpha using fund-specific benchmarks rather than the total market portfolio; benchmark selection introduces measurement error that allows aggregate alpha to differ from zero. Second, the Fama--French multifactor alphas strip out systematic factor exposures (size, value, profitability, investment), so aggregate alpha will be non-zero whenever active managers as a group carry systematic style tilts relative to the market. Third, our sample is a subset of the active management industry (U.S.\ open-end mutual funds and ETFs), not the complete universe of active investors, so the accounting identity binds only approximately.

These deviations are important because they mean that our dependent variable is not trivially zero by construction. Nevertheless, the Sharpe arithmetic implies that the scope for passive share to explain variation in aggregate gross alpha is inherently limited. The low $R^2$ values we observe in our regressions (1--6\%) are consistent with this constraint, and we discuss the implications in Section~\ref{sec:discussion}.

\subsubsection{The Berk--Green Equilibrium}
\label{subsub:berk_green}

\textcite{berk2004} proposed an alternative equilibrium model that predicts zero net alpha regardless of market structure. In their framework, investors rationally allocate capital to managers with demonstrated skill, driving up each manager's AUM until the marginal fund's net alpha is competed away by diminishing returns to scale. The key insight is that the competitive mechanism operates through capital flows to individual managers, not through the aggregate ratio of informed to uninformed capital.

This model provides a distinct prediction for our setting. If the Berk--Green mechanism is the dominant force, then aggregate net alpha should be approximately zero at all times, and passive share should be irrelevant. The growth of passive investing merely reallocates capital across the active-passive boundary but does not affect the equilibrium in which rational flows drive each active fund's net alpha to zero. Finding no association between passive share and alpha would be consistent with the Berk--Green equilibrium, offering an important alternative explanation for a null result that is distinct from the Grossman--Stiglitz prediction.

\subsection{Empirical Evidence on Active Fund Performance}

\subsubsection{The Distribution of Fund Skill}

A large body of empirical work has characterized the cross-sectional distribution of mutual fund performance. \textcite{carhart1997} demonstrated that after controlling for the momentum factor, there is virtually no evidence of performance persistence among mutual funds. Funds that performed well in the past do not reliably continue to outperform, suggesting that much of what appears to be skill is attributable to factor exposures and luck.

\textcite{fama2010} extended this analysis using a bootstrap approach to separate skill from luck. They found that the distribution of estimated fund alphas is consistent with a population in which few, if any, managers possess genuine stock-picking ability. The vast majority of active funds produce net returns indistinguishable from what would be expected by chance alone.

\textcite{barras2010} addressed the multiple-testing problem inherent in evaluating thousands of funds simultaneously. Using a false discovery rate methodology, they estimated that approximately 75\% of funds have zero alpha, while the remainder is split between genuinely skilled and genuinely unskilled managers. Importantly, the proportion of skilled managers declined over their sample period, suggesting intensifying competition.

\textcite{berk2015} offered a more optimistic perspective, arguing that the value added by mutual fund managers (measured as gross alpha times AUM) is economically large and positive. Under this view, competition drives fees to the point where net alpha is approximately zero, but this reflects the efficient pricing of managerial talent rather than the absence of skill. The distinction between gross and net alpha is central to our analysis.

\subsubsection{Scale, Competition, and Decreasing Returns to Skill}

\textcite{pastor2015} documented that the aggregate skill of the active management industry has increased over time, but that this improvement has been more than offset by the growth of the industry itself. As more capital chases alpha opportunities, the marginal return to skill declines. \textcite{pastor2012} formalized this in a model where the equilibrium size of the active management industry is determined by the trade-off between decreasing returns to scale at the industry level and investors' demand for active management.

These models, together with the \textcite{berk2004} framework, predict that competition erodes alpha through capital allocation rather than through the proportion of informed traders. If the competitive mechanism operates primarily through AUM flows, then the growth of passive investing may not increase the alpha available to active managers; it merely removes capital from the active pool without reducing the competitive intensity among the remaining participants.

\subsection{The Impact of Passive Investing on Market Efficiency}

\subsubsection{Evidence for Reduced Price Informativeness}

Several studies have documented that the growth of passive investing is associated with reduced informational efficiency at the individual security level. \textcite{israeli2017} found that higher ETF ownership is associated with wider bid-ask spreads, higher return synchronicity, and lower future earnings response coefficients, suggesting that prices become less informationally efficient as passive ownership increases.

\textcite{glosten2021} examined how ETF activity affects the informational efficiency of underlying securities and found mixed results. While ETF arbitrage can improve short-term price efficiency by correcting transient mispricing, the mechanical trading associated with index rebalancing can also introduce non-fundamental price pressure.

\textcite{qin2015} provided evidence that indexing leads to a deterioration in price informativeness, measured by the amount of firm-specific information reflected in stock prices. They argued that passive funds' mechanical purchasing of all stocks in an index, regardless of fundamentals, introduces correlated trading pressure that reduces the signal-to-noise ratio in prices.

\textcite{ben-david2018} documented that higher ETF ownership is associated with increased intraday volatility and non-fundamental price movements in the underlying securities, consistent with the hypothesis that passive investing introduces noise into prices.

\subsubsection{Counter-Evidence: Structural Impediments to Active Performance}

\textcite{busse2022} provided a critical counter-narrative to the Grossman--Stiglitz prediction. They documented that the unprecedented movement toward passive funds during the 2010s coincided with a challenging stock-picking environment created by the dominance of mega-capitalization stocks. When a small number of very large companies drive the majority of index returns, active managers who typically underweight their largest positions face a structural headwind.

Crucially, \textcite{busse2022} found that traditional performance predictors such as Active Share \parencite{cremers2009} and tracking error only predict future returns when alpha opportunities are plentiful. During periods of mega-cap dominance, even the most active managers significantly underperform, and performance fails to persist. This suggests that the relationship between passive growth and active performance may be more nuanced than the Grossman--Stiglitz framework implies: while reduced information gathering should theoretically create opportunities, market structure effects may independently constrain alpha generation.

\subsubsection{Liquidity, Crowding, and Correlated Trading Channels}

A separate strand of literature examines how the growth of passive investing affects the trading environment for active managers. \textcite{sushko2018} argued that passive funds' commitment to holding index constituents reduces the pool of available counterparties for active managers seeking to trade, particularly during periods of market stress. If active managers are increasingly trading with each other rather than with a diverse set of counterparties, their strategies may become more correlated, reducing the aggregate alpha available.

\textcite{coles2022} documented that index-linked investing creates correlated demand shocks across index constituents, making it harder for active managers to exploit firm-specific information. This ``crowding'' hypothesis is distinct from the Grossman--Stiglitz mechanism: rather than passive growth creating more mispricing, it may instead make the remaining active strategies more correlated and thus more vulnerable to common shocks.

The liquidity implications of this crowding are amplified during periods of market stress. \textcite{hameed2010} showed that stock market declines are associated with significant reductions in market liquidity, creating a feedback loop in which falling prices trigger further liquidity withdrawal. \textcite{brunnermeier2009} documented how such liquidity spirals were central to the 2007--2008 financial crisis, with leveraged institutions facing simultaneous margin calls that forced correlated asset sales. If the growth of passive investing reduces the diversity of active counterparties, these liquidity-spiral dynamics may become more severe during stress periods, as the remaining active managers hold increasingly correlated positions that unwind simultaneously. This theoretical channel motivates our exploration of stress-conditional interaction effects in the empirical analysis.

\subsection{Hypotheses}

Based on the theoretical and empirical literature reviewed above, we formulate the following hypotheses:

\begin{description}[leftmargin=2cm, style=nextline]
\item[H1 (Grossman--Stiglitz):] There is a positive association between the market share of passive investing and the aggregate, asset-weighted alpha of active fund managers.
\item[H2 (Null / Berk--Green):] There is no statistically significant association between passive market share and aggregate active alpha. This is consistent with the view that competitive forces \parencite{berk2004, pastor2012, pastor2015} and structural impediments \parencite{busse2022} offset any informational inefficiency created by passive growth.
\item[H3 (Conditional):] The association between passive share and alpha varies with market regime. During periods of market stress, the relationship may differ from the unconditional average, though the direction is ambiguous a priori: the Grossman--Stiglitz framework predicts a \textit{positive} interaction (more passive + more stress = more mispricing = more alpha), while the crowding and liquidity channels \parencite{sushko2018, coles2022, hameed2010} predict a \textit{negative} interaction (more passive + more stress = more correlated selling = less alpha).
\end{description}

We note that H3 is formulated as a two-sided hypothesis. The direction of the stress-conditional effect is a key empirical question rather than a pre-specified prediction.


% ============================================================================
% 3. DATA AND METHODOLOGY
% ============================================================================
\section{Data and Methodology}
\label{sec:methodology}

\subsection{Data Sources}
\label{subsec:data_sources}

\subsubsection{Fund Data}

The primary dataset is sourced from Morningstar Direct and covers the period from December 1996 to January 2025. The raw sample comprises approximately 29,800 actively managed and 3,700 passively managed U.S.\ open-end mutual funds and ETFs. For each fund, we obtain monthly gross returns, net returns, Net Asset Values (NAV), assets under management, expense ratios, prospectus benchmark identifiers, and Morningstar category classifications.

Critically, the Morningstar database includes both currently active funds and ``dead'' funds (those that were liquidated or merged during the sample period). This survivorship-bias-free design is essential for our analysis because funds that perform poorly are disproportionately likely to be liquidated, and excluding them would bias alpha estimates upward \parencite{elton1996}. We note, however, that we do not quantify the survivorship bias adjustment in our sample; the claim rests on the inclusion of dead funds, not on a formal decomposition.

\subsubsection{Factor Data}

Risk factor data for the Fama--French three-factor and five-factor models are obtained from Kenneth French's data library \parencite{french_data}. The factors include the market risk premium (Mkt$-$RF), size (SMB), value (HML), profitability (RMW), and investment (CMA), along with the risk-free rate (one-month Treasury bill rate). All factor returns are expressed in monthly percentage terms.

\subsubsection{Macroeconomic Controls}

We obtain macroeconomic control variables from the Federal Reserve Economic Data (FRED) database via its API. These include the CBOE Volatility Index (VIX), investment-grade credit spreads (the difference between Moody's Baa-rated corporate bond yield and the 10-year Treasury yield), and the market excess return (from the Fama--French factors). These variables serve as proxies for market uncertainty, financial stress, and aggregate market conditions, respectively.

We note a potential ``bad controls'' concern \parencite{angrist2009}: if passive investing affects alpha partly \textit{through} its influence on market volatility, credit conditions, or aggregate returns, then conditioning on these variables absorbs a portion of the mediating mechanism. We discuss this issue further in Section~\ref{subsec:econometrics} and present specifications both with and without these controls.

\subsection{Data Processing}
\label{subsec:data_processing}

\subsubsection{Deduplication}

Mutual funds frequently offer multiple share classes (e.g., Class A, Class C, Institutional) that correspond to the same underlying portfolio but differ in fee structures. Including multiple share classes would artificially inflate the sample size and double-count AUM. We identify and remove share classes labeled ``Load Waived,'' which are duplicates that exist solely for fee-waiver purposes. This procedure removes approximately 2,000 redundant entries, ensuring that each fund in our sample represents a unique investment portfolio.

\subsubsection{Stale Pricing Detection}

Illiquid funds or those with reporting errors may report identical NAV values for multiple consecutive periods, creating artificial stretches of zero volatility that would bias risk estimates downward. We implement a detection algorithm that flags ``NAV runs,'' defined as instances where the NAV remains unchanged for more than seven consecutive periods. Observations within flagged NAV runs are excluded from the alpha estimation.

\subsubsection{Benchmark Matching}

Accurate alpha estimation requires an appropriate benchmark for each fund. We implement a hierarchical matching procedure:
\begin{enumerate}
    \item \textbf{Primary match:} Each fund is matched to its designated prospectus benchmark as reported in fund documents.
    \item \textbf{Fallback match:} When the prospectus benchmark is unavailable or the corresponding index return data cannot be located (approximately 12.9\% of funds), we fall back to the Morningstar Category Index, which represents the average return of all funds in the same investment category.
\end{enumerate}
This hierarchical approach ensures that performance is measured against a relevant market proxy while maximizing sample coverage. In our final sample, 85.1\% of active funds are matched to their prospectus benchmark, 12.9\% to a Morningstar Category Index, and 2.0\% remain unmatched and are excluded from the alpha calculation.

\subsection{Variable Construction}
\label{subsec:variables}

\subsubsection{Dependent Variable: Risk-Adjusted Alpha}

We estimate monthly alpha for each active fund $i$ using three asset pricing models. Alpha is estimated via 36-month rolling ordinary least squares (OLS) regressions, requiring a minimum of 18 valid monthly observations within each estimation window \parencite{ferson1996}. The three specifications are:

\paragraph{Capital Asset Pricing Model (CAPM).} For the CAPM, we estimate a benchmark-relative Jensen's alpha. Rather than using the aggregate market portfolio as the single factor, we use each fund's matched benchmark as the reference:
\begin{equation}
    \alpha_{i} = R_{i,t} - \bigl[R_{f,t} + \hat{\beta}_{i}(R_{bench,i,t} - R_{f,t})\bigr]
    \label{eq:capm}
\end{equation}
where $R_{i,t}$ is the return of fund $i$ in month $t$, $R_{f,t}$ is the risk-free rate, $R_{bench,i,t}$ is the return of fund $i$'s matched benchmark, and $\hat{\beta}_{i}$ is estimated from the rolling OLS regression of fund excess returns on benchmark excess returns. This differs from the textbook CAPM, which uses the market portfolio as the single factor. Our benchmark-relative formulation measures skill relative to each fund's designated benchmark rather than absolute risk-adjusted performance. We adopt this approach because funds with different mandates (e.g., small-cap value, international equity) face different opportunity sets, and measuring alpha relative to a mandate-specific benchmark provides a more meaningful measure of stock-picking ability.

\paragraph{Fama--French Three-Factor Model (FF3).} This specification adds controls for size and value exposures using the standard Fama--French factors:
\begin{equation}
    R_{i,t} - R_{f,t} = \alpha_{i} + \beta_{mkt,i}\text{MKT}_t + \beta_{smb,i}\text{SMB}_t + \beta_{hml,i}\text{HML}_t + \varepsilon_{i,t}
    \label{eq:ff3}
\end{equation}
where MKT is the market excess return, and SMB and HML are the size and value factors \parencite{fama1993}. Unlike our CAPM specification, the FF3 and FF5 models use the standard Fama--French market factor rather than fund-specific benchmarks. This means the CAPM and multifactor alphas measure conceptually different objects: the CAPM alpha captures benchmark-relative skill, while the multifactor alphas capture returns unexplained by common risk factors. This distinction is relevant for interpreting differences in results across models.

\paragraph{Fama--French Five-Factor Model (FF5).} The most comprehensive specification adds controls for profitability and investment:
\begin{equation}
    R_{i,t} - R_{f,t} = \alpha_{i} + \beta_{mkt,i}\text{MKT}_t + \beta_{smb,i}\text{SMB}_t + \beta_{hml,i}\text{HML}_t + \beta_{rmw,i}\text{RMW}_t + \beta_{cma,i}\text{CMA}_t + \varepsilon_{i,t}
    \label{eq:ff5}
\end{equation}
where RMW and CMA are the profitability and investment factors \parencite{fama2015}. By controlling for the broadest set of systematic risk exposures, FF5 alpha isolates the component of returns that cannot be attributed to known factor tilts. We designate the five-factor model as our primary specification because it provides the most comprehensive risk adjustment.

\subsubsection{Gross vs.\ Net Alpha}
\label{subsub:gross_net}

A central methodological decision concerns the treatment of management fees. Net alpha, computed from after-fee returns, conflates managerial skill with fee decisions: $\alpha_{net} = \alpha_{gross} - \text{fees}$. Because our analysis includes fee-related variables as potential explanatory factors, using net alpha as the dependent variable introduces a mechanical relationship between the dependent variable and fee-based regressors. Specifically, regressing net alpha on fee spread (the difference between active and passive fees) is partially regressing fees on fees, a tautological relationship that can inflate $R^2$ and distort coefficient estimates.

To break this mechanical link, our primary analysis uses \textbf{gross alpha}, estimated from pre-fee returns. Gross alpha measures managerial stock-picking ability independent of the fee charged. We compute gross returns as:
\begin{equation}
    R_{gross,i,t} = R_{net,i,t} + \frac{\text{ExpenseRatio}_{i}}{12}
    \label{eq:gross}
\end{equation}
We retain net alpha specifications as a sensitivity check to demonstrate the exact mechanism by which fee controls distort the results (see Section~\ref{subsec:panel_c}).

\subsubsection{Aggregation}

Individual fund alphas are aggregated into a market-level time series by computing the asset-weighted average across all active funds in each month $t$:
\begin{equation}
    \bar{\alpha}_t = \frac{\sum_{i=1}^{N_t} \alpha_{i,t} \cdot \text{AUM}_{i,t}}{\sum_{i=1}^{N_t} \text{AUM}_{i,t}}
    \label{eq:aggregation}
\end{equation}
Asset-weighting is essential to reflect the economic reality of the market. It prevents small, volatile funds from disproportionately influencing the aggregate measure and represents the dollar-weighted experience of the average investor \parencite{french2008}. We note that this approach uses contemporaneous AUM as weights; using lagged AUM ($\text{AUM}_{i,t-1}$) would avoid a subtle look-ahead bias arising from the correlation between current-month returns and end-of-month AUM. We retain contemporaneous weighting for consistency with the prior literature, noting that the bias is small (one month's return has a marginal effect on total AUM).

The aggregation from approximately 30,000 fund-month observations to a single market-level time series of 302 monthly observations entails a substantial loss of cross-sectional information. This design choice is motivated by our research question, which concerns the market-level association between passive share and aggregate alpha, but it limits statistical power and prevents us from examining fund-level heterogeneity. We discuss the power implications in Section~\ref{subsec:power}.

\subsubsection{Primary Independent Variable: Passive Market Share}

The primary independent variable is the passive market share, computed monthly as:
\begin{equation}
    \text{PassiveShare}_t = \frac{\sum_{j \in \text{Passive}} \text{AUM}_{j,t}}{\sum_{j \in \text{Passive}} \text{AUM}_{j,t} + \sum_{i \in \text{Active}} \text{AUM}_{i,t}}
    \label{eq:passive_share}
\end{equation}
This variable captures the evolving structure of the asset management industry. In our sample, passive share increases from approximately 5.8\% in December 1996 to over 51\% by January 2025.

\subsubsection{Control Variables}

Our baseline control set includes three macroeconomic variables chosen to capture time-varying market conditions that may independently affect active alpha:

\begin{itemize}[noitemsep]
    \item \textbf{Market excess return} (Mkt$-$RF from Fama--French): Controls for the well-documented tendency of active managers to underperform in strong bull markets \parencite{fama2010}.
    \item \textbf{VIX}: The CBOE Volatility Index proxies for market uncertainty and expected volatility. Higher volatility may create more opportunities for active managers to exploit transient mispricings.
    \item \textbf{Investment-grade credit spread}: The spread between Moody's Baa corporate bond yield and the 10-year Treasury yield proxies for aggregate financial stress and credit risk appetite.
\end{itemize}

In the extended specification, we add:
\begin{itemize}[noitemsep]
    \item \textbf{Factor dispersion}: Defined as $|\text{SMB}| + |\text{HML}|$ for the CAPM and FF3 models, and $|\text{SMB}| + |\text{HML}| + |\text{RMW}| + |\text{CMA}|$ for the FF5 model. This variable captures the dispersion in returns across style factors. Higher factor dispersion indicates a market where style bets are particularly rewarded or penalized, creating more cross-sectional variation and potentially more opportunities for active managers. We use absolute values because large factor moves in either direction increase the cross-sectional dispersion of stock returns, expanding the opportunity set for stock pickers. Unlike alpha dispersion (which is endogenous), factor dispersion is constructed entirely from exogenous Fama--French factors.
\end{itemize}

In the sensitivity panel, we additionally include:
\begin{itemize}[noitemsep]
    \item \textbf{Fee spread}: The difference between asset-weighted active fund fees and passive fund fees. This variable is excluded from the main specification due to non-stationarity (ADF $p = 0.97$; KPSS rejects stationarity at 1\%) and its collinearity with passive share (both variables trend monotonically over the sample). It is included only in the sensitivity analysis (Panel C) to demonstrate how its inclusion distorts the estimated relationship.
\end{itemize}

We acknowledge the potential ``bad controls'' problem noted above \parencite{angrist2009}. If passive share affects alpha partly through its influence on market volatility (VIX), credit conditions, or aggregate market returns, then conditioning on these variables absorbs the mediating channel and attenuates the estimated coefficient on passive share. We address this concern in two ways. First, we report the baseline specification without controls (Panel A, specification 1) alongside the specification with controls (specification 2), allowing the reader to compare the coefficient magnitudes. Second, the first-differences specification (Panel B) omits the level controls by construction (it differences all variables), providing an alternative estimate that does not condition on potentially endogenous levels.

\subsection{Econometric Framework}
\label{subsec:econometrics}

\subsubsection{Main Specification}

The core time-series regression takes the form:
\begin{equation}
    \bar{\alpha}_{t}^{gross} = \beta_0 + \beta_1 \text{PassiveShare}_t + \boldsymbol{\beta}' \mathbf{X}_t + \varepsilon_t
    \label{eq:main_regression}
\end{equation}
where $\bar{\alpha}_{t}^{gross}$ is the asset-weighted aggregate gross alpha in month $t$, $\text{PassiveShare}_t$ is the passive market share, and $\mathbf{X}_t$ is a vector of control variables.

The coefficient of primary interest is $\beta_1$. Under the Grossman--Stiglitz hypothesis (H1), we expect $\beta_1 > 0$: higher passive share should be associated with higher alpha. Under the null hypothesis (H2), $\beta_1 = 0$.

\subsubsection{Model Hierarchy and Decision Rule}
\label{subsub:model_hierarchy}

We estimate all specifications under three factor models (CAPM, FF3, FF5), which differ in their risk adjustment and, for the CAPM, in their benchmark definition. To guard against selective model emphasis, we adopt the following decision rule: we designate FF5 as our primary specification because it provides the most comprehensive risk adjustment, and we consider a finding robust only if it achieves statistical significance ($p < 0.05$) in at least two of the three factor models. Where results differ across models, we report all three and discuss the discrepancy rather than selectively emphasizing the most favorable outcome.

\subsubsection{Standard Errors}

Time-series regressions of aggregate portfolio returns are susceptible to serial correlation due to the overlapping nature of the 36-month rolling estimation windows used to construct fund-level alphas. This rolling window induces a moving-average structure of order up to 35 in the alpha residuals.

To obtain valid inference, we employ Newey--West heteroskedasticity and autocorrelation consistent (HAC) standard errors \parencite{newey1987} with a Bartlett kernel. Our baseline lag truncation parameter is $L = 6$, selected using the standard rule $L = \lfloor 4(T/100)^{2/9} \rfloor$ for $T \approx 300$. However, this data-driven rule does not account for the known MA structure induced by the rolling window, so $L = 6$ may underestimate the true autocorrelation and produce standard errors that are too small. We address this concern by reporting results for $L = 6$, $L = 12$, and $L = 36$ in the robustness panel. The choice of $L = 36$ matches the rolling window length and provides the most conservative inference. If the CAPM result (which is closest to the significance boundary at $L = 6$) loses significance at $L = 36$, we note this prominently.

\subsubsection{Stationarity}
\label{subsub:stationarity}

Passive market share exhibits a strong upward trend over our sample period and is clearly non-stationary. Augmented Dickey--Fuller (ADF) tests \parencite{dickey1979} fail to reject the unit root null ($p > 0.95$), and KPSS tests \parencite{kwiatkowski1992} reject the stationarity null ($p < 0.01$). The dependent variable, gross alpha, is approximately stationary (ADF rejects the unit root null). Regressing a stationary I(0) variable on a non-stationary I(1) regressor can produce spurious inference \parencite{phillips1986}.

The mixed integration orders of our variables mean that standard cointegration frameworks (Engle--Granger, Johansen) do not apply, since cointegration requires all variables to share the same integration order. We address the non-stationarity concern in two ways. First, we include macroeconomic control variables that absorb some of the time-series variation. Second, and more importantly, we estimate a first-differences specification as a co-primary analysis:
\begin{equation}
    \Delta \bar{\alpha}_{t}^{gross} = \gamma_0 + \gamma_1 \Delta \text{PassiveShare}_t + \boldsymbol{\gamma}' \Delta \mathbf{X}_t + \eta_t
    \label{eq:first_diff}
\end{equation}
which eliminates deterministic and stochastic trends, albeit at the cost of focusing on short-run dynamics rather than long-run relationships. We treat the first-differences specification as co-primary alongside the levels specification, not merely as a robustness check. Readers who are skeptical of the levels regression (on the grounds that OLS on mixed I(0)/I(1) series has non-standard inference) should weight the first-differences results more heavily.

\subsubsection{Structural Break Test}
\label{subsub:structural_break}

The sensitivity of our results to the inclusion of the dot-com period (1999--2002), in which the FF3 coefficient changes sign, motivates a formal test for structural instability. We implement a Chow test \parencite{chow1960} at the boundary between the dot-com period and the subsequent sample (December 2002). A significant Chow statistic would indicate that the relationship between passive share and alpha differs between the two sub-periods, justifying split-sample analysis or the use of regime-specific controls.

\subsubsection{Diagnostics}

For each specification, we report the Durbin--Watson statistic \parencite{durbin1951} as a preliminary indicator of residual autocorrelation, the Breusch--Godfrey test \parencite{breusch1978, godfrey1978} with 30 lags for a formal serial correlation test, and condition numbers to assess multicollinearity among regressors.

\subsubsection{Power Analysis}
\label{subsec:power}

The aggregate time-series design compresses approximately 30,000 fund-month observations into 302 market-level data points. This raises the question of whether our test has sufficient statistical power to detect economically meaningful effects.

We provide a back-of-the-envelope power calculation. The standard deviation of monthly gross alpha is approximately 0.6--0.8\% across models, and the standard deviation of passive share over the sample is approximately 0.14 (14 percentage points). For an effect size of $\beta_1 = -0.5$ (a 10-percentage-point increase in passive share associated with a $-0.05$\% monthly decline in alpha, or $-0.6$\% annualized), the approximate $t$-statistic under the alternative is:
\begin{equation}
    t \approx \frac{|\beta_1| \cdot \text{SD}(\text{PS}) \cdot \sqrt{N}}{\text{SD}(\alpha)} = \frac{0.5 \times 0.14 \times \sqrt{302}}{0.7} \approx 1.74
    \label{eq:power}
\end{equation}
This corresponds to approximately 50\% power at the 5\% significance level, using one-sided critical values. Power at the conventional 80\% threshold would require either a larger sample, a stronger effect, or less noise. This calculation underscores that the FF3 and FF5 null results cannot be interpreted as ``no effect''; they are equally consistent with ``effect too small to detect.'' We return to this limitation in the Discussion.

\subsubsection{Multiple Testing Correction for Exploratory Analysis}
\label{subsub:multiple_testing}

Panel D of our results reports five exploratory specifications (lagged passive share at two horizons, quadratic, and two interaction terms), each estimated across three factor models, yielding 15 hypothesis tests. With a conventional 5\% significance threshold and 15 tests, we would expect approximately one rejection by chance alone. To guard against spurious findings, we apply the Benjamini--Hochberg false discovery rate (FDR) procedure \parencite{benjamini1995} to the Panel D $p$-values and report both raw and FDR-adjusted $p$-values. We treat Panel D results as hypothesis-generating regardless of their statistical significance.

\subsection{Specification Overview}

Our analysis is organized into four panels, each addressing a distinct aspect of the passive share--alpha association. Table~\ref{tab:spec_overview} provides an overview.

\begin{table}[H]
\centering
\caption{Overview of regression specifications}
\label{tab:spec_overview}
\small
\begin{tabularx}{\textwidth}{llX}
\toprule
\textbf{Panel} & \textbf{Specifications} & \textbf{Purpose} \\
\midrule
A: Main Analysis & (1) Baseline, (2) Controls, (3) Extended, (4) No Dot-Com & Estimate the core association using gross alpha \\
B: Robustness & (5) First-Differences, (6--8) NW lag=3/12/36 & Address stationarity and inference sensitivity \\
C: Sensitivity & (8) Gross+fee, (9) Net no-fee, (10) Net+fee & Demonstrate the fee-spread sign-flip mechanism \\
D: Exploratory & (11--12) Lagged PS, (13) Quadratic, (14--15) Interactions & Test predictive, non-linear, and conditional effects \\
\bottomrule
\end{tabularx}
\end{table}


% ============================================================================
% 4. RESULTS
% ============================================================================
\section{Results}
\label{sec:results}

\subsection{Descriptive Statistics}
\label{subsec:descriptives}

\subsubsection{Passive Share Growth}

Figure~\ref{fig:passive_share} traces the evolution of passive market share over our sample period. Passive share grew from 5.8\% in December 1996 to 51.2\% by January 2025, crossing the 50\% threshold in August 2024. Growth was gradual through the early 2000s, accelerated following the 2008 financial crisis, and continued at a rapid pace through the 2010s and 2020s.

\begin{figure}[H]
    \centering
    \includegraphics[width=\textwidth]{figures/passive_share_timeline.pdf}
    \caption{Evolution of passive market share in the U.S.\ mutual fund industry. Passive share grew from 5.8\% in December 1996 to 51.2\% in January 2025, crossing the 50\% threshold in August 2024. The horizontal dashed line marks the 50\% level. Key market events and NBER recession periods (grey shading) are annotated. Data source: Morningstar Direct.}
    \label{fig:passive_share}
\end{figure}

\subsubsection{Fee Compression}

Concurrent with the growth of passive investing, the asset management industry experienced significant fee compression. Asset-weighted active fund fees declined from approximately 1.0\% per year to 0.65\%, while passive fund fees fell from approximately 0.25\% to below 0.10\%. Figure~\ref{fig:fee_analysis} illustrates these trends.

\begin{figure}[H]
    \centering
    \includegraphics[width=\textwidth]{figures/fee_analysis.pdf}
    \caption{Fee dynamics in the U.S.\ mutual fund industry (1996--2025). Top left: asset-weighted average fees for active and passive funds. Top right: the active-passive fee differential. Bottom left: simple vs.\ asset-weighted average fees for active funds. Bottom right: fund count over time by type. Data source: Morningstar Direct.}
    \label{fig:fee_analysis}
\end{figure}

\subsubsection{Aggregate Alpha}

Table~\ref{tab:alpha_descriptive} reports descriptive statistics for aggregate alpha across the three asset pricing models.

\begin{table}[H]
\centering
\caption{Descriptive statistics for monthly aggregate asset-weighted gross alpha (\%)}
\label{tab:alpha_descriptive}
\small
\begin{tabular}{lrrrrr}
\toprule
\textbf{Model} & \textbf{Mean} & \textbf{Median} & \textbf{Std.\ Dev.} & \textbf{Min} & \textbf{Max} \\
\midrule
CAPM   & 0.096 & 0.079 & 0.610 & $-$3.064 & 5.313 \\
FF3    & 0.147 & 0.109 & 0.777 & $-$2.208 & 2.859 \\
FF5    & 0.209 & 0.143 & 0.784 & $-$1.999 & 2.784 \\
\bottomrule
\end{tabular}
\begin{flushleft}
\small\textit{Note.} Statistics computed over the 302 months with valid alpha estimates (December 1999 to January 2025). The 36-month rolling estimation window requires data from December 1996, meaning the first valid alpha observation is December 1999. Alpha is expressed in monthly percentage terms.
\end{flushleft}
\end{table}

All three models produce mean gross alphas close to zero, consistent with the Sharpe (1991) accounting identity discussed in Section~\ref{subsub:sharpe_arithmetic}. CAPM alpha is closest to zero (0.096\% per month), while FF3 and FF5 alphas are somewhat higher (0.147\% and 0.209\%, respectively), reflecting the fact that the multifactor models attribute more return variation to systematic factor exposures, leaving a larger residual when active managers carry systematic style tilts.

Figure~\ref{fig:rolling_alpha} plots the aggregate gross alpha over time. The time series reveals substantial volatility in the early period (1999--2003), corresponding to the dot-com bubble, followed by a gradual convergence toward zero. The 12-month moving average highlights the declining trend across all models.

\begin{figure}[H]
    \centering
    \includegraphics[width=\textwidth]{figures/rolling_alpha_all_models.pdf}
    \caption{Aggregate asset-weighted gross alpha over time for CAPM (top), Fama--French 3-factor (middle), and Fama--French 5-factor (bottom). Thin lines show monthly values; dashed lines show 12-month moving averages. Gray shaded areas mark NBER recession periods.}
    \label{fig:rolling_alpha}
\end{figure}

\subsubsection{Period Heterogeneity}

Table~\ref{tab:alpha_periods} decomposes the alpha statistics by sub-period. The dot-com era (1999--2002) stands out with notably higher alpha across all models (CAPM: 0.265\%, FF3: 0.670\%, FF5: 0.707\% monthly), reflecting the exceptional volatility and return dispersion of the technology bubble. This period combines very low passive share ($<$15\%) with high alpha, creating a strong anchor point that influences full-sample regression estimates.

\begin{table}[H]
\centering
\caption{Aggregate alpha by sub-period (monthly \%)}
\label{tab:alpha_periods}
\small
\begin{tabular}{lrrrrrr}
\toprule
\textbf{Period} & \multicolumn{2}{c}{\textbf{CAPM}} & \multicolumn{2}{c}{\textbf{FF3}} & \multicolumn{2}{c}{\textbf{FF5}} \\
\cmidrule(lr){2-3} \cmidrule(lr){4-5} \cmidrule(lr){6-7}
 & Mean & Std & Mean & Std & Mean & Std \\
\midrule
Dot-com (1999--2002)    & 0.265 & 1.324 & 0.670 & 1.169 & 0.707 & 1.214 \\
Pre-GFC (2003--2007)    & 0.162 & 0.400 & $-$0.145 & 0.673 & $-$0.045 & 0.634 \\
Recovery (2008--2014)   & 0.070 & 0.472 &  0.239 & 0.784 &  0.321 & 0.826 \\
Bull (2015--2019)       & 0.047 & 0.346 &  0.087 & 0.522 &  0.129 & 0.505 \\
Recent (2020--2024)     & 0.000 & 0.436 &  0.044 & 0.609 &  0.089 & 0.593 \\
\bottomrule
\end{tabular}
\begin{flushleft}
\small\textit{Note.} Monthly gross alpha in percentage points. N months: Dot-com 37, Pre-GFC 60, Recovery 84, Bull 60, Recent 61.
\end{flushleft}
\end{table}

Figure~\ref{fig:period_comparison} provides a visual comparison of mean gross alpha across sub-periods and models. The dot-com era stands out clearly, with all three models showing their highest alpha. The subsequent decline is monotonic for CAPM but more nuanced for the multifactor models, which show a secondary peak during the Recovery period (2008--2014).

\begin{figure}[H]
    \centering
    \includegraphics[width=\textwidth]{figures/period_comparison.pdf}
    \caption{Mean monthly gross alpha by sub-period and asset pricing model. The dot-com period (December 1999 to December 2002) exhibits the highest alpha across all models. The general decline in alpha coincides with the rising passive market share.}
    \label{fig:period_comparison}
\end{figure}

\subsection{Stationarity Tests}

Table~\ref{tab:stationarity} reports ADF and KPSS test results for the key variables. Gross alpha, market excess return, VIX, and credit spread are confirmed stationary. Passive share and fee spread are non-stationary, warranting caution in levels regressions and motivating the first-differences co-primary specification.

\begin{table}[H]
\centering
\caption{Stationarity tests for regression variables}
\label{tab:stationarity}
\small
\begin{tabular}{lrrrrl}
\toprule
\textbf{Variable} & \textbf{ADF stat} & \textbf{ADF $p$} & \textbf{KPSS stat} & \textbf{KPSS $p$} & \textbf{Verdict} \\
\midrule
Gross alpha (CAPM)   & $-$15.32 & $<$0.001 & 0.257 & $>$0.10 & Stationary \\
Passive share        & 3.22     & 1.000    & 2.849 & $<$0.01 & Non-stationary \\
Market excess return & $-$17.10 & $<$0.001 & 0.101 & $>$0.10 & Stationary \\
VIX                  & $-$4.70  & $<$0.001 & 0.354 & 0.097   & Stationary \\
Credit spread (IG)   & $-$3.61  & 0.006    & 0.214 & $>$0.10 & Stationary \\
Factor dispersion    & $-$3.03  & 0.032    & 0.383 & 0.085   & Stationary \\
Fee spread           & 0.24     & 0.974    & 2.625 & $<$0.01 & Non-stationary \\
\bottomrule
\end{tabular}
\begin{flushleft}
\small\textit{Note.} ADF: Augmented Dickey--Fuller test; H$_0$: unit root. KPSS: Kwiatkowski--Phillips--Schmidt--Shin test; H$_0$: stationarity. Both tests agree that gross alpha (CAPM) is stationary (ADF rejects the unit root; KPSS does not reject stationarity).
\end{flushleft}
\end{table}

\subsection{Panel A: Main Analysis}
\label{subsec:panel_a}

Table~\ref{tab:panel_a} reports the main regression results. We estimate the association between passive share and gross alpha across four specifications for each of the three alpha models.

\begin{table}[H]
\centering
\caption{Panel A: Gross alpha regressed on passive share (Newey--West HAC)}
\label{tab:panel_a}
\small
\begin{tabular}{lrrrr}
\toprule
 & \textbf{(1) Baseline} & \textbf{(2) Controls} & \textbf{(3) Extended} & \textbf{(4) No Dot-Com} \\
\midrule
\multicolumn{5}{l}{\textit{Panel A.1: CAPM}} \\
PassiveShare    & $-$0.597**\phantom{*}  & $-$0.620**\phantom{*}  & $-$0.607**\phantom{*}  & $-$0.346\phantom{***} \\
                & (0.278)               & (0.293)               & (0.267)               & (0.212) \\
Mkt$-$RF       &                       & 0.027***              & 0.027***              & 0.023*** \\
                &                       & (0.008)               & (0.008)               & (0.007) \\
VIX             &                       & $-$0.001\phantom{***} & $-$0.005\phantom{***} & $-$0.003\phantom{***} \\
                &                       & (0.008)               & (0.006)               & (0.005) \\
Credit spread   &                       & 0.035\phantom{***}    & 0.064\phantom{***}    & 0.069\phantom{***} \\
                &                       & (0.089)               & (0.074)               & (0.067) \\
Factor disp.    &                       &                       & 0.021\phantom{***}    &  \\
                &                       &                       & (0.028)               &  \\[6pt]
$N$             & 302  & 302  & 302  & 265  \\
$R^2$           & 0.016 & 0.058 & 0.065 & 0.084 \\
DW              & 1.993 & 2.054 & 2.067 & 1.825 \\
\midrule
\multicolumn{5}{l}{\textit{Panel A.2: FF3}} \\
PassiveShare    & $-$0.528\phantom{***}  & $-$0.433\phantom{***}  & $-$0.404\phantom{***}  & 0.574**\phantom{*} \\
                & (0.521)               & (0.553)               & (0.464)               & (0.276) \\
$R^2$           & 0.008 & 0.015 & 0.038 & 0.037 \\
\midrule
\multicolumn{5}{l}{\textit{Panel A.3: FF5 (Primary)}} \\
PassiveShare    & $-$0.709\phantom{***}  & $-$0.652\phantom{***}  & $-$0.603\phantom{***}  & 0.234\phantom{***} \\
                & (0.504)               & (0.539)               & (0.448)               & (0.297) \\
$R^2$           & 0.014 & 0.016 & 0.037 & 0.020 \\
\bottomrule
\end{tabular}
\begin{flushleft}
\small\textit{Note.} Dependent variable: monthly aggregate asset-weighted gross alpha (\%). Standard errors in parentheses are Newey--West HAC with 6 lags. Significance: *** $p<0.01$, ** $p<0.05$, * $p<0.10$. Controls in specifications (2)--(4): market excess return, VIX, investment-grade credit spread. Specification (3) adds factor dispersion. Specification (4) excludes the dot-com period (1999--2002). FF3 and FF5 panels show the passive share coefficient only; control variable estimates are qualitatively similar to the CAPM panel.
\end{flushleft}
\end{table}

\textbf{CAPM.} The passive share coefficient is negative and statistically significant across specifications (1)--(3), with the controls specification yielding $\beta_1 = -0.620$ ($p = 0.035$). The economic magnitude is modest: a 10 percentage point increase in passive share is associated with a $-0.062$\% per month ($-0.74$\% annualized) decline in gross CAPM alpha. When the dot-com period is excluded (specification 4), the coefficient weakens to $-0.346$ ($p = 0.102$), indicating that the 1999--2002 period contributes substantially to the full-sample result.

\textbf{FF3 and FF5.} Under our primary specification (FF5), the passive share coefficient is $-0.652$ ($p = 0.227$), which is negative but statistically insignificant. The FF3 coefficient is similarly insignificant ($-0.433$, $p = 0.434$). These multifactor models attribute more return variation to style factors, leaving less unexplained variation and reducing the power to detect an association with passive share. We note that the Breusch--Godfrey test rejects the null of no serial correlation for both multifactor models (FF3: $p = 0.014$; FF5: $p = 0.006$), while it does not reject for the CAPM ($p = 0.94$). The Newey--West standard errors are designed to be robust to such autocorrelation, but the rejection confirms that the rolling window induces more residual persistence in the multifactor alpha series. Notably, the FF3 coefficient reverses sign to positive and significant when the dot-com period is excluded ($+0.574$, $p = 0.037$), though this result is not confirmed by FF5 ($+0.234$, $p = 0.431$) and is therefore not robust by our cross-model criterion.

Applying our decision rule (Section~\ref{subsub:model_hierarchy}), the negative association between passive share and gross alpha is significant only in the CAPM and does not achieve significance in the FF3 or FF5 models. By our stated criterion of requiring significance in at least two of three models, the main effect of passive share is \textit{not robust}. Figure~\ref{fig:forest} summarizes the passive share coefficients across all Panel A specifications and models.

\begin{figure}[H]
    \centering
    \includegraphics[width=\textwidth]{figures/coefficient_forest.pdf}
    \caption{Forest plot of passive share coefficients across Panel A specifications and factor models. Points show coefficient estimates; horizontal bars show 95\% confidence intervals based on Newey--West HAC standard errors (6 lags). The vertical dashed line marks zero. Only the CAPM estimates exclude zero from their confidence intervals.}
    \label{fig:forest}
\end{figure}

Figure~\ref{fig:scatter_all_models} provides a visual representation of the bivariate relationship. The scatter plots reveal the high dispersion of monthly alpha and the modest slope of the fitted line.

\begin{figure}[H]
    \centering
    \includegraphics[width=\textwidth]{figures/scatter_all_models.pdf}
    \caption{Scatter plots of monthly gross alpha against passive market share for CAPM (left), FF3 (center), and FF5 (right). Points are color-coded by time (purple = early, yellow = recent). Red lines show OLS fits.}
    \label{fig:scatter_all_models}
\end{figure}

\subsection{Panel B: Robustness and First-Differences}
\label{subsec:panel_b}

\subsubsection{First-Differences (Co-Primary Specification)}

As discussed in Section~\ref{subsub:stationarity}, the first-differences specification addresses the non-stationarity of passive share and serves as a co-primary analysis alongside the levels regression. Table~\ref{tab:panel_b} reports the results.

\begin{table}[H]
\centering
\caption{Panel B: First-differences and Newey--West lag sensitivity}
\label{tab:panel_b}
\small
\begin{tabular}{lrrr}
\toprule
\textbf{Specification} & \textbf{CAPM} & \textbf{FF3} & \textbf{FF5} \\
\midrule
\multicolumn{4}{l}{\textit{First-differences ($\Delta$PassiveShare):}} \\
$\gamma_1$ & $-$209.5*** & $-$42.3 & $-$34.9 \\
$p$-value  & 0.001       & 0.285   & 0.383 \\
$R^2$      & 0.150       & 0.022   & 0.016 \\[6pt]
\multicolumn{4}{l}{\textit{NW lag sensitivity (Controls specification):}} \\
Lag $= 3$:  $p$-value & 0.046  & 0.359 & 0.162 \\
Lag $= 6$:  $p$-value & 0.035  & 0.434 & 0.227 \\
Lag $= 12$: $p$-value & 0.038  & 0.517 & 0.319 \\
Lag $= 24$: $p$-value & 0.027  & 0.571 & 0.380 \\
Lag $= 36$: $p$-value & 0.020  & 0.583 & 0.393 \\
\bottomrule
\end{tabular}
\begin{flushleft}
\small\textit{Note.} First-differences regresses $\Delta$alpha on $\Delta$passive\_share and $\Delta$controls. NW lag sensitivity uses the controls specification (Panel A, spec 2) with varying Newey--West lag truncation parameters. The coefficient estimate ($\beta_1 = -0.620$) is identical across NW lags; only the standard errors and $p$-values change. *** $p<0.01$, ** $p<0.05$, * $p<0.10$.
\end{flushleft}
\end{table}

The first-differences CAPM coefficient is $\gamma_1 = -209.5$ ($p = 0.001$), confirming the negative association in a specification that eliminates trends. The large magnitude reflects the very small month-to-month changes in passive share: $\Delta$PassiveShare has a standard deviation of approximately 0.0015 (0.15 percentage points). To facilitate comparison with the levels estimate, we compute the standardized coefficient. A one-standard-deviation increase in $\Delta$PassiveShare is associated with a $-209.5 \times 0.0015 \approx -0.31$\% decline in monthly alpha, or roughly 0.5 standard deviations of the dependent variable. In contrast, the levels estimate implies a 10-percentage-point increase in passive share (a change that unfolds over many years) is associated with a $-0.062$\% monthly decline. The two specifications are therefore answering different questions: the levels regression captures the long-run average gradient, while first-differences captures month-to-month co-movement. The sign agreement is reassuring, though the orders-of-magnitude difference in raw coefficients warrants the standardized comparison.

For FF3 and FF5, the first-differences coefficients are negative but statistically insignificant, consistent with the levels results.

\subsubsection{Newey--West Lag Sensitivity}

A key concern raised in Section~\ref{subsec:econometrics} is that the 36-month rolling window induces MA(35) autocorrelation, potentially making our baseline NW lag of $L = 6$ too short. Table~\ref{tab:panel_b} reports $p$-values for $L \in \{3, 6, 12, 24, 36\}$.

The CAPM result is robust to, and in fact \textit{strengthened by}, longer lag truncation: the $p$-value decreases from 0.035 at $L = 6$ to 0.020 at $L = 36$, the most conservative choice matching the rolling window length. This addresses the concern that the CAPM significance might be an artifact of underestimated standard errors. The FF3 and FF5 $p$-values increase with longer lags but remain far from any conventional significance threshold throughout. Figure~\ref{fig:nw_sensitivity} visualizes this pattern.

\begin{figure}[H]
    \centering
    \includegraphics[width=\textwidth]{figures/nw_lag_sensitivity.pdf}
    \caption{Newey--West lag sensitivity for the passive share coefficient across all three models. The CAPM $p$-value (left axis) decreases with longer lags, confirming robustness. FF3 and FF5 remain insignificant throughout.}
    \label{fig:nw_sensitivity}
\end{figure}

\subsection{Panel C: The Fee-Spread Artifact}
\label{subsec:panel_c}

Panel C presents our methodological contribution. In an earlier stage of this research, we employed a specification that included fee spread (the difference between asset-weighted active and passive fund fees) as a control variable. This produced a positive and significant coefficient on passive share, apparently supporting the Grossman--Stiglitz hypothesis. We subsequently discovered that this result is an artifact of fee spread's properties. Table~\ref{tab:panel_c} demonstrates the mechanism.

\begin{table}[H]
\centering
\caption{Panel C: Fee-spread sensitivity, the sign-flip mechanism}
\label{tab:panel_c}
\small
\begin{tabular}{llrrr}
\toprule
\textbf{Spec.} & \textbf{Description} & \textbf{CAPM $\beta_{PS}$} & \textbf{FF3 $\beta_{PS}$} & \textbf{FF5 $\beta_{PS}$} \\
\midrule
(2) & Gross $\alpha$, controls              & $-$0.620**   & $-$0.433     & $-$0.652     \\
(8) & Gross $\alpha$ + fee\_spread          & $+$4.256*    & $+$4.926***  & $+$3.638**   \\
(9) & Net $\alpha$, controls (no fee)       & $-$0.532*    & $-$0.341     & $-$0.558     \\
(10)& Net $\alpha$ + fee\_spread            & $+$4.381*    & $+$4.956***  & $+$3.661**   \\
\bottomrule
\end{tabular}
\begin{flushleft}
\small\textit{Note.} All specifications include market excess return, VIX, and credit spread as controls. Newey--West HAC standard errors with 6 lags. *** $p<0.01$, ** $p<0.05$, * $p<0.10$. The key observation is that the sign flip occurs in \textit{both} gross and net alpha specifications with nearly identical magnitudes.
\end{flushleft}
\end{table}

Figure~\ref{fig:panel_c} visualizes the sign flip. The sign of the passive share coefficient flips from negative to positive when fee spread is added, in all three models and in both gross and net alpha specifications.

\begin{figure}[H]
    \centering
    \includegraphics[width=\textwidth]{figures/panel_c_sign_flip.pdf}
    \caption{The fee-spread sign-flip mechanism. Each pair of bars shows the passive share coefficient without (left) and with (right) fee spread as a control variable. The sign reversal is consistent across all three factor models and across gross and net alpha specifications.}
    \label{fig:panel_c}
\end{figure} The near-identical magnitudes of the flip for gross alpha (specification 8: CAPM $+4.256$) and net alpha (specification 10: CAPM $+4.381$) are diagnostic. If the sign flip were driven by the mechanical fee component in net alpha ($\alpha_{net} = \alpha_{gross} - \text{fees}$), it should be substantially larger for net alpha than gross alpha. Instead, the coefficients are virtually indistinguishable, confirming that the distortion arises from the collinearity between fee spread and passive share, both of which trend monotonically over the sample, rather than from any mechanical fee contamination.

Fee spread fails the ADF unit root test ($p = 0.974$) and has a condition number of approximately 13,000 when included alongside passive share, versus approximately 180 without it. Adding a non-stationary, collinear regressor to a regression on a stationary dependent variable absorbs the trend component of passive share, leaving only the residual variation, which happens to be positively correlated with alpha in the multifactor models.

We emphasize that we initially employed this specification ourselves and present Panel C not as a critique of other researchers, but as documentation of a methodological pitfall that may tempt future researchers working with fee and performance data. The lesson is general: any study regressing net fund performance on fee-related controls risks conflating the mechanical fee relationship with the economic relationship of interest.

\subsection{Panel D: Exploratory Analysis}
\label{subsec:panel_d}

Panel D examines predictive, non-linear, and stress-conditional extensions. We present these as exploratory analyses and apply the Benjamini--Hochberg false discovery rate (FDR) correction \parencite{benjamini1995} to account for the multiple comparisons involved. Table~\ref{tab:panel_d} reports both raw and FDR-adjusted $p$-values.

\begin{table}[H]
\centering
\caption{Panel D: Exploratory specifications with multiple-testing correction}
\label{tab:panel_d}
\small
\begin{tabular}{llrrrrrr}
\toprule
 & & \multicolumn{2}{c}{\textbf{CAPM}} & \multicolumn{2}{c}{\textbf{FF3}} & \multicolumn{2}{c}{\textbf{FF5}} \\
\cmidrule(lr){3-4} \cmidrule(lr){5-6} \cmidrule(lr){7-8}
\textbf{Spec.} & \textbf{Variable} & $\beta$ & $p$ / $p_{BH}$ & $\beta$ & $p$ / $p_{BH}$ & $\beta$ & $p$ / $p_{BH}$ \\
\midrule
(11) & PS lag 3m     & $-$0.598 & .040 / .082 & $-$0.419 & .449 / .483 & $-$0.645 & .231 / .388 \\
(12) & PS lag 12m    & $-$0.629 & .041 / .082 & $-$0.392 & .483 / .483 & $-$0.631 & .246 / .388 \\
(13) & PS (linear)   & $-$3.440 & .022 / .082 & $-$4.357 & .261 / .418 & $-$3.481 & .346 / .388 \\
(13) & PS$^2$        & $+$4.932 & .037 / .082 & $+$6.863 & .255 / .418 & $+$4.948 & .388 / .388 \\
(14) & PS $\times$ VIX     & $+$0.021 & .553 / .632 & $-$0.122 & .034 / .137 & $-$0.053 & .357 / .388 \\
(15) & PS $\times$ Credit  & $+$0.203 & .701 / .701 & $-$2.201 & .016 / .125 & $-$1.655 & .058 / .388 \\
\bottomrule
\end{tabular}
\begin{flushleft}
\small\textit{Note.} All specifications include market excess return, VIX, and credit spread as controls. NW HAC standard errors with 6 lags. Raw $p$-values are reported before the slash; Benjamini--Hochberg FDR-adjusted $p$-values after the slash. No result survives the FDR correction at the 5\% level.
\end{flushleft}
\end{table}

\subsubsection{Lagged Passive Share}

Specifications (11) and (12) replace contemporaneous passive share with 3-month and 12-month lags, respectively. CAPM coefficients are $-0.598$ ($p = 0.040$) and $-0.629$ ($p = 0.041$), directionally consistent with the levels result. However, these regressions provide limited information about temporal precedence, because passive share is a near-unit-root process (lag-1 autocorrelation $>0.99$). Lagged passive share is essentially identical to contemporaneous passive share, so the lagged regressions do not meaningfully address reverse causality. After FDR correction, neither result is significant ($p_{BH} = 0.082$).

\subsubsection{Non-Linearity}

Specification (13) adds a quadratic term. For the CAPM, both the linear ($-3.440$, $p = 0.022$) and quadratic ($+4.932$, $p = 0.037$) terms are significant at raw thresholds, suggesting a U-shaped relationship with a turning point at approximately 35\% passive share. However, neither term survives FDR correction ($p_{BH} = 0.082$), and the result does not replicate in FF3 or FF5.

\subsubsection{Stress-Conditional Interactions}

The interaction between passive share and credit spread is the most interesting raw result: FF3 yields $\beta_{PS \times Credit} = -2.201$ ($p = 0.016$), with borderline FF5 support ($-1.655$, $p = 0.058$). The negative sign implies that during periods of wide credit spreads (financial stress), high passive concentration is associated with significantly lower multifactor alpha, the opposite of the Grossman--Stiglitz prediction. However, after FDR correction, the FF3 result is no longer significant ($p_{BH} = 0.125$), and the FF5 result is far from significance ($p_{BH} = 0.388$). The CAPM interaction is negligible ($+0.203$, $p = 0.701$). We therefore treat this finding as suggestive and hypothesis-generating, not as established evidence. Figure~\ref{fig:stress_contour} visualizes the interaction surface.

\begin{figure}[H]
    \centering
    \includegraphics[width=0.85\textwidth]{figures/stress_conditional_contour.pdf}
    \caption{Stress-conditional interaction between passive share and credit spread for FF3 alpha. The contour plot shows predicted alpha as a function of passive share (x-axis) and credit spread (y-axis). Darker regions indicate lower predicted alpha. The negative interaction implies that the combination of high passive share and high credit spread is associated with the lowest alpha.}
    \label{fig:stress_contour}
\end{figure}

\subsection{Panel E: Extended Diagnostics}
\label{subsec:panel_e}

\subsubsection{Power Analysis}

Table~\ref{tab:power} reports the minimum detectable effect size at 80\% power for each model. These estimates, derived from the residual standard error and sample size of the controls specification, indicate that our test is well-powered only for the CAPM.

\begin{table}[H]
\centering
\caption{Power analysis: minimum detectable $|\beta_1|$ at 80\% power ($\alpha = 0.05$)}
\label{tab:power}
\small
\begin{tabular}{lrrrr}
\toprule
\textbf{Model} & \textbf{Residual SE} & \textbf{Detectable $|\beta_1|$} & \textbf{Annual equiv.} & \textbf{Actual $|\hat{\beta}_1|$} \\
\midrule
CAPM & 0.596 & 0.825 & 9.9\%/yr & 0.620 \\
FF3  & 0.776 & 1.554 & 18.6\%/yr & 0.433 \\
FF5  & 0.782 & 1.516 & 18.2\%/yr & 0.652 \\
\bottomrule
\end{tabular}
\begin{flushleft}
\small\textit{Note.} The ``annual equivalent'' translates the monthly coefficient to the annualized change in alpha associated with a 100-percentage-point increase in passive share ($|\beta_1| \times 12$). The actual FF3 and FF5 coefficients are well below their respective detection thresholds, confirming that the null results for these models are uninformative about the true effect size.
\end{flushleft}
\end{table}

The CAPM's actual coefficient ($|\hat{\beta}_1| = 0.620$) is below its detection threshold (0.825), implying that even the CAPM result is borderline in terms of power. For FF3 and FF5, the detection thresholds (1.554 and 1.516) are two to three times larger than the estimated coefficients, confirming that these models lack sufficient power to detect economically plausible effects. The FF3 and FF5 null results cannot distinguish ``no association'' from ``association too small to detect with 302 observations.''

\subsubsection{Structural Break Test}

The Chow test \parencite{chow1960} at the March 2003 boundary (end of the dot-com period) rejects the null of parameter stability for all three models (CAPM: $F = 4.24$, $p < 0.001$; FF3: $F = 13.52$, $p < 0.001$; FF5: $F = 11.37$, $p < 0.001$). The pre-break sub-sample (December 1999 to February 2003, $N = 39$) yields very large negative passive share coefficients (CAPM: $-24.97$; FF3: $-110.81$; FF5: $-105.06$), driven by the coincidence of rapidly rising alpha and slowly rising passive share during the dot-com bubble. The post-break sub-sample (March 2003 to January 2025, $N = 263$) yields coefficients near zero (CAPM: $-0.37$; FF3: $+0.61$; FF5: $+0.23$).

This structural break confirms the outsized influence of the dot-com period and suggests that the full-sample negative association in the CAPM is driven primarily by the contrast between the pre-2003 and post-2003 regimes, rather than by a stable association throughout the sample.

\subsubsection{Zero-Alpha Threshold Estimates}

Following the original analysis, we compute the passive share level at which the baseline regression line crosses zero. However, we now report bootstrap confidence intervals (5,000 iterations, circular block bootstrap with block length 36) to convey the substantial uncertainty around these estimates. The block bootstrap preserves the temporal dependence structure in the regression residuals, providing more conservative confidence intervals than IID resampling.

\begin{table}[H]
\centering
\caption{Bootstrap confidence intervals for zero-alpha passive share thresholds}
\label{tab:zero_alpha_bootstrap}
\small
\begin{tabular}{lrrrrr}
\toprule
\textbf{Model} & \textbf{Point est.} & \textbf{95\% CI lower} & \textbf{95\% CI upper} & \textbf{$R^2$} & \textbf{\% neg.\ slope} \\
\midrule
CAPM & 43.0\% & 33.6\% &  71.1\% & 0.016 & 99.5\% \\
FF3  & 46.5\% & $-$78.5\% & 157.3\% & 0.008 & 73.7\% \\
FF5  & 49.0\% & $-$115.5\% & 184.5\% & 0.014 & 79.9\% \\
\bottomrule
\end{tabular}
\begin{flushleft}
\small\textit{Note.} Point estimates computed as $-\hat{\beta}_0 / \hat{\beta}_1$ from the baseline regression (specification 1). Bootstrap 95\% CIs from 5,000 iterations using a circular block bootstrap (block length 36) to preserve residual autocorrelation structure. The wide intervals reflect the low $R^2$ values: passive share explains 1--2\% of the variation in alpha, making the zero-crossing point highly uncertain.
\end{flushleft}
\end{table}

The bootstrap confidence intervals are extremely wide, particularly for the multifactor models. The CAPM interval spans from 33.6\% to 71.1\%, while the FF3 and FF5 intervals include negative values (FF3: $-78.5$\% to $157.3$\%; FF5: $-115.5$\% to $184.5$\%), meaning the regression line may not cross zero within any economically meaningful range. The block bootstrap, which accounts for residual serial correlation, yields a tighter CAPM interval than an IID bootstrap would, but substantially widens the FF3 and FF5 intervals and lowers the percentage of negative slopes (73.7\% and 79.9\% respectively, compared to 99.5\% for CAPM). These results underscore that the zero-alpha threshold is not a reliable estimate and should be interpreted with great caution. We report these estimates for completeness but do not build our conclusions upon them.

Figure~\ref{fig:residual_diagnostics} presents residual diagnostic plots for the CAPM controls specification, confirming that the Newey--West standard errors are appropriate for the observed residual structure.

\begin{figure}[H]
    \centering
    \includegraphics[width=\textwidth]{figures/residual_diagnostics.pdf}
    \caption{Residual diagnostics for the CAPM controls specification (Panel A, specification 2). Top left: residuals over time. Top right: residual histogram with normal overlay. Bottom left: residual autocorrelation function. Bottom right: Q-Q plot. The residuals exhibit modest leptokurtosis but no systematic pattern, consistent with the Breusch--Godfrey non-rejection ($p = 0.94$).}
    \label{fig:residual_diagnostics}
\end{figure}


% ============================================================================
% 5. DISCUSSION
% ============================================================================
\section{Discussion}
\label{sec:discussion}

\subsection{Summary of Findings}

Our analysis yields four main findings. First, under the CAPM, there is a statistically significant negative association between passive share and aggregate gross alpha ($\beta = -0.620$, $p = 0.035$), robust to Newey--West lag selection up to $L = 36$ ($p = 0.020$) and confirmed in first-differences ($p = 0.001$). Second, this result does not generalize to the Fama--French three-factor or five-factor models, which show negative but statistically insignificant coefficients. Under our pre-specified decision rule (significance in at least two of three models, with FF5 as primary), the main effect of passive share is not robust. Third, the inclusion of fee spread as a control variable reverses the sign of the passive share coefficient in all models, a specification artifact we document in Panel C. Fourth, exploratory interaction effects between passive share and credit spreads, while intriguing at raw significance levels (FF3: $p = 0.016$), do not survive the Benjamini--Hochberg FDR correction ($p_{BH} = 0.125$).

\subsection{Interpreting the CAPM Result}

The CAPM negative association is our most robust individual finding. It survives all lag lengths, the dot-com exclusion weakens it to marginal insignificance ($p = 0.102$), and the first-differences specification provides strong confirmation ($p = 0.001$). The NW lag sensitivity analysis is particularly reassuring: rather than losing significance as the lag increases (as reviewers of earlier drafts anticipated), the $p$-value \textit{decreases} from 0.035 at $L = 6$ to 0.020 at $L = 36$.

Taken at face value, this result is directionally \textit{opposite} to the Grossman--Stiglitz prediction (H1) and supports the view that passive growth is associated with lower, not higher, benchmark-relative alpha. It does not ``challenge the narrative that passive investing erodes active management alpha''; if anything, the CAPM result is consistent with that narrative.

However, we caution against overinterpreting a single-model finding. The CAPM alpha is benchmark-relative (Equation~\ref{eq:capm}), measuring skill against each fund's specific benchmark, while the multifactor alphas use a common set of risk factors. These measure conceptually different objects, and the discrepancy in results may reflect the different risk adjustments rather than a robust economic relationship. The Chow test confirms a structural break at 2003, with the pre-2003 sub-sample driving most of the full-sample association. The CAPM result is therefore best characterized as a statistically significant but model-specific and period-sensitive negative association.

\subsection{The FF3/FF5 Null Result and the Berk--Green Equilibrium}

Under our primary specification (FF5), we find no statistically significant association between passive share and aggregate gross alpha ($\beta = -0.652$, $p = 0.227$). However, the power analysis (Table~\ref{tab:power}) reveals that FF5 would require $|\beta_1| > 1.52$ for detection at 80\% power, more than twice the estimated coefficient. The null result is therefore uninformative: it cannot distinguish ``no association'' from ``association too small to detect.''

The null result is consistent with the \textcite{berk2004} equilibrium, in which rational capital allocation drives each active fund's net alpha to zero, regardless of the fraction of capital in passive vehicles. Under this view, the competitive mechanism operates through AUM flows to individual managers (investors chase performance, driving up AUM until alpha is dissipated), not through the aggregate ratio of informed to uninformed capital. The growth of passive investing merely reallocates capital across the active-passive boundary without affecting the Berk--Green equilibrium. Our finding that aggregate gross alpha fluctuates around a near-zero mean throughout the sample, even as passive share grew from 6\% to 51\%, is consistent with this framework.

It is equally consistent with the \textcite{pastor2012, pastor2015} models of decreasing returns to scale, in which the active management industry has grown more skilled over time but this improvement has been offset by increased competition. Under either interpretation, passive share is a spectator rather than a driver of aggregate alpha dynamics.

\subsection{The Fee-Spread Artifact}

Panel C documents what we consider our most robust methodological finding. The inclusion of fee spread as a control variable reverses the sign of the passive share coefficient from negative to positive in all three factor models and in both gross and net alpha specifications. The near-identical magnitude of the flip across gross and net alpha ($+4.256$ vs.\ $+4.381$ for CAPM) eliminates the mechanical fee channel as the explanation and points to a pure collinearity effect: both passive share and fee spread trend monotonically over the sample, so conditioning on one reverses the sign of the other.

We emphasize that we discovered this artifact in our own earlier analysis. We initially employed a specification regressing net alpha on passive share with fee-spread controls and obtained a positive, significant coefficient that appeared to support the Grossman--Stiglitz hypothesis. The correction to gross alpha, combined with the removal of fee spread, produced the null result that we now report. We present Panel C not as a critique of other researchers but as a cautionary example. Any study regressing fund performance (net or gross) on fee-related controls that share the same temporal trend risks the same distortion.

\subsection{Stress-Conditional Effects}

The FF3 interaction between passive share and credit spread ($\beta = -2.201$, raw $p = 0.016$) is directionally interesting: it suggests that during periods of financial stress, high passive concentration is associated with lower multifactor alpha. This is the opposite of the Grossman--Stiglitz prediction (which implies that stress and high passive share should \textit{increase} mispricing and thus alpha) and is directionally consistent with crowding or liquidity-withdrawal channels \parencite{sushko2018, coles2022, hameed2010}.

However, this finding does not survive the Benjamini--Hochberg FDR correction ($p_{BH} = 0.125$), is significant in only one of three factor models, and the CAPM interaction is negligible ($+0.203$, $p = 0.701$). We cannot identify the specific channel through which the interaction operates; the crowding interpretation is one of several plausible mechanisms, and distinguishing among them would require holdings-based measures of strategy similarity or fund-level flow data that are beyond the scope of this thesis. We treat this finding as hypothesis-generating: the negative interaction between passive share and financial stress merits further investigation with more powerful research designs.

\subsection{The Bad Controls Problem}

Our main specification controls for market excess return, VIX, and credit spread. These variables are included to absorb time-varying market conditions that may independently affect alpha. However, if passive investing affects alpha partly \textit{through} its influence on these variables, as some theoretical channels suggest (e.g., passive inflows may compress volatility or inflate market returns), then the controls absorb the mediating mechanism, attenuating the estimated coefficient on passive share \parencite{angrist2009}.

Two observations are relevant. First, the baseline specification without controls (Panel A, specification 1) yields a passive share coefficient of $-0.597$ ($p = 0.032$) for the CAPM, which is virtually identical to the controls specification ($-0.620$, $p = 0.035$). The controls absorb very little of the passive share effect, suggesting that mediation through these specific variables is limited. Second, the first-differences specification, which differences all variables and does not condition on levels, produces a strongly significant result ($p = 0.001$), further suggesting that the controls are not absorbing a large mediating channel. Nevertheless, we cannot rule out mediation through channels we do not measure, and we note this as a limitation.

\subsection{Limitations}

We identify several limitations that qualify our findings.

\textbf{Statistical power.} Our aggregate time-series design compresses approximately 30,000 fund-month observations into 302 market-level data points. The power analysis (Table~\ref{tab:power}) shows that our test has approximately 50\% power for the CAPM and substantially less for the multifactor models. The null results for FF3 and FF5 should not be interpreted as ``no association'' but rather as ``insufficient evidence to detect the association, if one exists.''

\textbf{Aggregation.} The aggregation to market-level alpha discards all fund-level heterogeneity. The Grossman--Stiglitz mechanism may operate at the cross-section (some managers benefit from reduced competition while others do not) rather than at the aggregate level. A fund-level panel regression with fund fixed effects would be more powerful but represents a substantially different research design.

\textbf{No causal identification.} Our design identifies an association, not a causal effect. We lack a credible instrument for passive share growth and cannot exploit a natural experiment. Reverse causality is a particular concern: low alpha may cause flows to passive (investors flee underperformance), generating a negative association that reflects the response of investors to poor active performance rather than any effect of passive share on alpha.

\textbf{The Sharpe arithmetic constraint.} As discussed in Section~\ref{subsub:sharpe_arithmetic}, the \textcite{sharpe1991} accounting identity constrains AUM-weighted aggregate gross alpha to fluctuate around approximately zero. Our measured alpha can deviate from zero due to benchmark selection, factor model specification, and incomplete coverage of the active management universe. Nevertheless, the constraint limits the scope for passive share to ``explain'' a variable that is pinned near zero by construction, and it contributes to the low $R^2$ values (1--6\%) observed in our regressions.

\textbf{Dot-com sensitivity and structural break.} The Chow test confirms a structural break at March 2003. The full-sample association is driven primarily by the pre-2003 period, in which the coincidence of low passive share and high alpha during the dot-com bubble creates a strong anchor. Post-2003, the association is near zero in all models. Whether the dot-com period represents a genuine manifestation of the passive-alpha relationship or a confounded episode (the bubble produced high alpha for reasons unrelated to passive share) is an open question.

\textbf{Stationarity.} Passive share is non-stationary (I(1)), while gross alpha is stationary (I(0)). Our levels regression therefore involves mixed integration orders, which can produce non-standard inference \parencite{phillips1986}. The first-differences specification addresses this concern, and we treat it as co-primary. Readers who are skeptical of the levels results should weight the first-differences estimates more heavily.

\textbf{Reverse causality.} The lagged passive share regressions (Panel D) do not meaningfully address reverse causality because passive share is a near-unit-root process. Lagged passive share is essentially identical to contemporaneous passive share (autocorrelation $>0.99$), so the lagged regressions add minimal information beyond what the contemporaneous regression provides. A proper Granger-causality test, including lagged alpha as a predictor, would be needed to address temporal precedence, and we leave this for future work.

\textbf{Benchmark matching.} Approximately 12.9\% of active funds are benchmarked against the Morningstar Category Index, which includes the fund's own return. For large funds that represent a substantial fraction of their category's AUM, this creates a circularity that mechanically attenuates measured alpha. We have not quantified this attenuation bias.


% ============================================================================
% 6. CONCLUSION
% ============================================================================
\section{Conclusion}
\label{sec:conclusion}

This thesis examines the association between the growth of passive investing and aggregate active management alpha in the U.S.\ mutual fund market, using a survivorship-bias-free sample of approximately 29,800 active and 3,700 passive funds from December 1996 to January 2025.

Our findings do not support the Grossman--Stiglitz hypothesis that passive growth creates alpha opportunities for remaining active managers. Under the CAPM, we find a statistically significant \textit{negative} association between passive share and gross alpha ($\beta = -0.620$, $p = 0.020$ with Newey--West lag 36), confirmed in first-differences ($p = 0.001$). However, this result does not generalize to the Fama--French three-factor or five-factor models, which yield negative but statistically insignificant coefficients. Under our pre-specified decision rule (FF5 as primary; robustness requires significance in at least two models), the main effect of passive share is not robust. The power analysis reveals that the multifactor null results are uninformative: our test lacks sufficient power to detect economically plausible effect sizes in these models. The finding is therefore best characterized as a CAPM-specific negative association that cannot be confirmed or rejected in the multifactor framework.

Our most robust contribution is methodological. We demonstrate that the inclusion of fee spread as a control variable reverses the sign of the passive share coefficient from negative to positive in all three factor models. This artifact arises from the collinearity of two trending variables, not from any mechanical fee contamination, and it operates identically in both gross and net alpha specifications. We document this pitfall as a cautionary example for researchers working with fee and performance data.

The exploratory analysis of stress-conditional effects yields an intriguing but fragile finding: the interaction between passive share and credit spread is negative and significant in the FF3 model at raw thresholds ($p = 0.016$), suggesting that high passive concentration during financial stress is associated with lower alpha. However, this result does not survive the Benjamini--Hochberg false discovery rate correction ($p_{BH} = 0.125$), is confined to a single factor model, and the underlying channel (crowding, liquidity withdrawal, or other mechanisms) cannot be identified with our data. We treat this as a hypothesis for future research.

Several directions for future work emerge from our analysis. First, a fund-level panel regression with fund fixed effects would exploit the cross-sectional heterogeneity that our aggregate design discards and would substantially improve statistical power. Second, the structural break at 2003 suggests that the passive-alpha association may be regime-dependent; a longer post-2020 sample, when passive share exceeds 50\%, would provide more observations in the high-passive-share regime where the Grossman--Stiglitz mechanism is most likely to bind. Third, the stress-conditional interaction between passive share and credit spreads, if confirmed with more powerful methods, would have implications for understanding systemic risk in an increasingly passive market. Fourth, holdings-based measures of strategy crowding and active share could provide direct evidence on the mechanisms that our aggregate analysis can only speculate about.

For researchers, our findings highlight the importance of specification choices in studies of passive investing and active performance. The fee-spread artifact demonstrates how a trending control variable can reverse the sign of the primary coefficient. For practitioners and asset allocators, our results offer neither support for the claim that passive growth creates opportunities for active managers nor strong evidence that it destroys them. The association, to the extent that it exists, is model-dependent, period-sensitive, and too small to be reliably detected with existing aggregate data.


% ============================================================================
% REFERENCES
% ============================================================================
\newpage
\printbibliography[title={References}]


% ============================================================================
% APPENDIX
% ============================================================================
\newpage
\appendix
\section{Appendix}
\label{sec:appendix}

\subsection{Full Regression Output: CAPM Controls Specification}
\label{app:full_output}

Table~\ref{tab:app_full_capm} reports the complete OLS output for the CAPM controls specification (Panel A, specification 2), which serves as the basis for most of the discussion.

\begin{table}[H]
\centering
\caption{Full regression output: CAPM gross alpha, controls specification}
\label{tab:app_full_capm}
\small
\begin{tabular}{lrrrr}
\toprule
\textbf{Variable} & \textbf{Coef.} & \textbf{NW SE} & \textbf{$p$-value} & \textbf{95\% CI} \\
\midrule
Constant              & 0.210  & 0.114 & 0.065 & [$-$0.013, 0.433] \\
PassiveShare          & $-$0.620 & 0.293 & 0.035 & [$-$1.195, $-$0.045] \\
Market excess return  & 0.027  & 0.008 & $<$0.001 & [0.012, 0.042] \\
VIX                   & $-$0.001 & 0.008 & 0.905 & [$-$0.016, 0.014] \\
Credit spread (IG)    & 0.035  & 0.089 & 0.694 & [$-$0.139, 0.209] \\
\midrule
$N$ & \multicolumn{4}{l}{302} \\
$R^2$ / Adj.\ $R^2$ & \multicolumn{4}{l}{0.058 / 0.045} \\
Durbin--Watson & \multicolumn{4}{l}{2.054} \\
Breusch--Godfrey $p$ (30 lags) & \multicolumn{4}{l}{0.944} \\
Condition number & \multicolumn{4}{l}{181} \\
\bottomrule
\end{tabular}
\begin{flushleft}
\small\textit{Note.} Newey--West HAC standard errors with 6 lags. The Breusch--Godfrey test does not reject the null of no serial correlation at any conventional level. The condition number of 181 indicates no multicollinearity concern.
\end{flushleft}
\end{table}

\subsection{Newey--West Lag Sensitivity: Full Results}
\label{app:nw_sensitivity}

Table~\ref{tab:app_nw} extends the Panel B lag sensitivity analysis with additional detail for all three models.

\begin{table}[H]
\centering
\caption{Newey--West lag sensitivity for passive share coefficient (controls specification)}
\label{tab:app_nw}
\small
\begin{tabular}{lrrrrrr}
\toprule
 & \multicolumn{2}{c}{\textbf{CAPM}} & \multicolumn{2}{c}{\textbf{FF3}} & \multicolumn{2}{c}{\textbf{FF5}} \\
\cmidrule(lr){2-3} \cmidrule(lr){4-5} \cmidrule(lr){6-7}
\textbf{NW lag} & SE & $p$ & SE & $p$ & SE & $p$ \\
\midrule
3  & 0.310 & 0.046 & 0.472 & 0.359 & 0.467 & 0.162 \\
6  & 0.293 & 0.035 & 0.553 & 0.434 & 0.539 & 0.227 \\
12 & 0.298 & 0.038 & 0.668 & 0.517 & 0.654 & 0.319 \\
24 & 0.280 & 0.027 & 0.764 & 0.571 & 0.743 & 0.380 \\
36 & 0.266 & 0.020 & 0.788 & 0.583 & 0.763 & 0.393 \\
\bottomrule
\end{tabular}
\begin{flushleft}
\small\textit{Note.} The coefficient ($\beta_1$) is identical across all lags (CAPM: $-0.620$; FF3: $-0.433$; FF5: $-0.652$). Only the standard errors and $p$-values change. The CAPM result strengthens with longer lags; FF3 and FF5 weaken but remain insignificant throughout.
\end{flushleft}
\end{table}

\subsection{Panel D: Benjamini--Hochberg FDR Correction Detail}
\label{app:fdr_detail}

Table~\ref{tab:app_fdr} reports the full set of raw and FDR-adjusted $p$-values for the key variables in Panel D across all three models. The correction is applied separately within each model.

\begin{table}[H]
\centering
\caption{BH-FDR corrected $p$-values for Panel D exploratory specifications}
\label{tab:app_fdr}
\small
\begin{tabular}{llrrr}
\toprule
\textbf{Model} & \textbf{Variable} & \textbf{Coef.} & \textbf{Raw $p$} & \textbf{$p_{BH}$} \\
\midrule
\multicolumn{5}{l}{\textit{CAPM}} \\
 & PS lag 3m        & $-$0.598 & 0.040 & 0.082 \\
 & PS lag 12m       & $-$0.629 & 0.041 & 0.082 \\
 & PS (linear, quad.)& $-$3.440 & 0.022 & 0.082 \\
 & PS$^2$           & $+$4.932 & 0.037 & 0.082 \\
 & PS $\times$ VIX  & $+$0.021 & 0.553 & 0.632 \\
 & PS $\times$ Credit & $+$0.203 & 0.701 & 0.701 \\
\midrule
\multicolumn{5}{l}{\textit{FF3}} \\
 & PS lag 3m        & $-$0.419 & 0.449 & 0.483 \\
 & PS lag 12m       & $-$0.392 & 0.483 & 0.483 \\
 & PS (linear, quad.)& $-$4.357 & 0.261 & 0.418 \\
 & PS$^2$           & $+$6.863 & 0.255 & 0.418 \\
 & PS $\times$ VIX  & $-$0.122 & 0.034 & 0.137 \\
 & PS $\times$ Credit & $-$2.201 & 0.016 & 0.125 \\
\midrule
\multicolumn{5}{l}{\textit{FF5}} \\
 & PS lag 3m        & $-$0.645 & 0.231 & 0.388 \\
 & PS lag 12m       & $-$0.631 & 0.246 & 0.388 \\
 & PS (linear, quad.)& $-$3.481 & 0.346 & 0.388 \\
 & PS$^2$           & $+$4.948 & 0.388 & 0.388 \\
 & PS $\times$ VIX  & $-$0.053 & 0.357 & 0.388 \\
 & PS $\times$ Credit & $-$1.655 & 0.058 & 0.388 \\
\bottomrule
\end{tabular}
\begin{flushleft}
\small\textit{Note.} Benjamini--Hochberg correction applied separately within each model across all Panel D key variables. No result survives the FDR correction at the 5\% level. The FF3 PS $\times$ Credit interaction (raw $p = 0.016$, $p_{BH} = 0.125$) would require a raw $p < 0.008$ to survive the correction.
\end{flushleft}
\end{table}

\subsection{Structural Break: Pre- and Post-2003 Coefficients}
\label{app:chow}

Table~\ref{tab:app_chow} reports the Chow test results and sub-sample coefficients for the structural break at March 2003.

\begin{table}[H]
\centering
\caption{Chow structural break test at March 2003}
\label{tab:app_chow}
\small
\begin{tabular}{lrrrrr}
\toprule
\textbf{Model} & \textbf{$F$-stat} & \textbf{$p$-value} & \textbf{$\beta_{PS}$ pre-2003} & \textbf{$\beta_{PS}$ post-2003} & \textbf{$N$ pre / post} \\
\midrule
CAPM & 4.24  & $<$0.001 & $-$24.97  & $-$0.37 & 39 / 263 \\
FF3  & 13.52 & $<$0.001 & $-$110.81 & $+$0.61 & 39 / 263 \\
FF5  & 11.37 & $<$0.001 & $-$105.06 & $+$0.23 & 39 / 263 \\
\bottomrule
\end{tabular}
\begin{flushleft}
\small\textit{Note.} Chow test for parameter equality between December 1999 to February 2003 and March 2003 to January 2025. All models strongly reject the null of parameter stability. The extreme pre-2003 coefficients reflect the small sample and the coincidence of low passive share with high dot-com alpha.
\end{flushleft}
\end{table}


\end{document}
